\section{The SBD Foundation: TransNetV2 and AutoShot}

In this section, we first formulate the shot boundary detection (SBD) problem. Next, we detail the spatial-temporal factorized convolution, dilated convolutional blocks, and similarity features that form the backbone of modern SBD architectures. Finally, we review the Neural Architecture Search (NAS) strategy employed by AutoShot to select the optimal backbone.

\subsection{Problem Formulation}

Given a video represented as an ordered sequence of $T$ frames $\{x_t\}_{t=1}^{T}$, SBD aims to predict a binary sequence $\{b_t\}_{t=1}^{T}$, where $b_t=1$ denotes a shot boundary (either an abrupt cut or a gradual transition) at frame $t$, and $b_t=0$ represents a non-boundary frame. An SBD model outputs a boundary probability sequence $\{p_t\}_{t=1}^{T}$ where $p_t \in [0,1]$. The final boundary set is obtained by thresholding the probability scores and optionally grouping consecutive positive predictions:
\begin{equation}
    \hat{b}_t =
    \begin{cases}
        1, & p_t > \theta,\\
        0, & p_t \leq \theta,
    \end{cases}
\end{equation}
where the threshold $\theta$ is optimized on a validation set to balance precision and recall.

\subsection{Spatial-Temporal Factorized Convolution}

Deep learning models for SBD must process spatial-temporal representations directly from raw frame sequences. To avoid the high parameter count and computational complexity of standard 3-D convolutional kernels of size $k \times k \times k$, modern architectures factorize each 3-D convolution into two sequential operations:
\begin{enumerate}
    \item \textbf{A 2-D spatial convolution} with kernel size $1 \times k \times k$. This processes each frame in the sequence independently, learning visual representation cues such as edges, textures, and color regions.
    \item \textbf{A 1-D temporal convolution} with kernel size $k \times 1 \times 1$. This processes the spatial features along the temporal axis to capture motion patterns and scene changes.
\end{enumerate}

With input/output channels $C$ and kernel size $k=3$, factorizing the kernel reduces the parameter count from $27C^2$ to $9C^2 + 3C^2 = 12C^2$. This represents a parameter reduction of approximately $56\%$, which mitigates overfitting and imposes an inductive bias that separates spatial contents from temporal motion.

\subsection{Dilated DCNN Blocks}

The backbone of the baseline architecture consists of stacked Dilated Convolutional Neural Network (DDCNN) blocks. Each block splits the input features into four parallel temporal convolution branches with dilation rates $d$ increasing exponentially:
\begin{equation}
    d \in \{1, 2, 4, 8\}.
\end{equation}
The outputs of these four branches are concatenated along the channel dimension, followed by Batch Normalization (BatchNorm) and Rectified Linear Unit (ReLU) activation.

This multi-scale dilation design expands the temporal receptive field. For a temporal kernel size $K=3$, the cumulative receptive field $R_k$ after block $k$ is calculated inductively as:
\begin{equation}
    R_k = R_{k-1} + 2 \cdot (K - 1) \cdot \sum_{d \in \{1,2,4,8\}} d = R_{k-1} + 60,
\end{equation}
with $R_0 = 1$ frame. Stacking six blocks theoretically yields a receptive field of up to $361$ frames. In practice, videos are processed in sliding windows of $100$ frames, so the effective temporal context is capped by the window length rather than the theoretical $361$ frames. This context is wide enough to capture long gradual transitions, such as dissolves and fades.

\subsection{Similarity Branch}

To supplement the features learned by the convolutional blocks, a parallel similarity branch is integrated into the model. This branch extracts two forms of visual consistency:
\begin{enumerate}
    \item \textbf{RGB Histogram Similarity}: A 512-bin color histogram is computed for each frame. The cosine distance between histograms of consecutive frames yields a 101-dimensional vector representing color shifts.
    \item \textbf{Learned Similarity}: The cosine similarity between normalized deep features extracted from the intermediate layers of the convolutional blocks is calculated.
\end{enumerate}
The histogram and learned similarities are fused and processed through a fully connected layer with 128 units and ReLU activation. The resulting similarity features are concatenated with the convolutional features before final classification.

\subsection{AutoShot Backbone Search}

While TransNetV2 employs a fixed configuration of identical DDCNN blocks, AutoShot \cite{zhu2023autoshot} optimizes the backbone configuration for short-form videos using Neural Architecture Search (NAS). The search space covers four DDCNN block variants (differing in channel configurations and spatial feature sharing) across 6 layers, and a choice of 0 to 4 Transformer layers \cite{vaswani2017attention} at the 7th layer, resulting in $83.9 \times 10^6$ candidate architectures. 

The search utilizes a Single-Path One-Shot supernetwork \cite{guo2020spos} combined with Bayesian Optimization. The baseline supernetwork is also compared against differentiable architectures like DARTS \cite{liu2019darts}. After searching, the model is further refined via knowledge distillation \cite{hinton2015distilling} and weight grafting. The optimal architecture found under the F1-score criterion (AutoShot@F1) consists of:
\begin{itemize}
    \item Layer 1: DDCNNV2 ($4$ filters, dilation $nd=4$).
    \item Layers 2--4: DDCNNV2A ($4$ filters, dilation $nd=5$) shared spatial convolutions.
    \item Layer 5: DDCNNV2 ($12$ filters, dilation $nd=5$).
    \item Layer 6: DDCNNV2 ($8$ filters, dilation $nd=5$).
    \item Layer 7: $0$ Transformer layers.
\end{itemize}
This backbone extracts a 4864-dimensional feature vector at the \texttt{fc1\_0} stage. In this paper, we freeze this optimized AutoShot backbone and denote the extracted features for frame $t$ as:
\begin{equation}
    h_t \in \mathbb{R}^{4864}.
\end{equation}

\section{Proposed AutoShotV2: Calibrated SBD with Dual-Branch Head}

AutoShotV2 improves the decision stage of SBD by training a new classification head and adding probability calibration and temporal smoothing post-processing steps. Fig. \ref{fig:pipeline} illustrates the overall inference pipeline.

\begin{figure*}[!t]
\centering
\resizebox{\textwidth}{!}{%
\begin{tikzpicture}[
    node distance=0.65cm and 0.9cm,
    font=\footnotesize,
    block/.style={draw, rounded corners=2pt, thick, align=center, minimum height=0.9cm, text width=2.45cm},
    wide/.style={draw, rounded corners=2pt, thick, align=center, minimum height=0.9cm, text width=3.0cm},
    arr/.style={-{Stealth[length=5pt]}, thick}
]
\node[block, fill=blue!10] (in) {Video frames\\$48\times27$ RGB};
\node[wide, fill=green!12, right=of in] (backbone) {Frozen AutoShot backbone\\3-D DCNN + similarity + histogram};
\node[block, fill=green!8, right=of backbone] (feat) {Cached features\\4864-D};
\node[block, fill=orange!15, right=of feat] (head) {New classification head\\Focal Loss};
\node[block, fill=purple!12, right=of head] (temp) {Temperature scaling\\$T^*=\PaperTemperature$};
\node[block, fill=pink!14, right=of temp] (smooth) {Gaussian smoothing\\$\sigma=\PaperGaussianSigma$};
\node[block, fill=blue!8, right=of smooth] (thr) {Threshold\\$\theta=\PaperDeployThreshold$};
\node[block, fill=green!18, right=of thr] (out) {Shot boundaries\\F1 = \PaperASVTwoShotFOne};
 
\draw[arr] (in) -- (backbone);
\draw[arr] (backbone) -- (feat);
\draw[arr] (feat) -- (head);
\draw[arr] (head) -- (temp);
\draw[arr] (temp) -- (smooth);
\draw[arr] (smooth) -- (thr);
\draw[arr] (thr) -- (out);
\end{tikzpicture}
}
\caption{AutoShotV2 inference pipeline. The searched AutoShot backbone is frozen. AutoShotV2 modifies the decision stage by training a new classification head and adding calibration and temporal smoothing before thresholding.}
\label{fig:pipeline}
\end{figure*}

\subsection{Network Architecture of the Classification Head}

The proposed head takes the cached feature vector $h_t$ and maps it to a shared hidden state:
\begin{equation}
    u_t = \mathrm{ReLU}(W_f h_t + c_f),
\end{equation}
where $W_f \in \mathbb{R}^{1024\times4864}$. A dropout layer with a keep probability of $0.5$ is applied to $u_t$ for regularization. Two parallel linear classification layers then output the logits:
\begin{equation}
    z_t^{(1)} = W_1 u_t + c_1,\qquad
    z_t^{(2)} = W_2 u_t + c_2.
\end{equation}
The primary branch ($z_t^{(1)}$) is trained with the exact one-hot boundary labels and is utilized during inference. The auxiliary branch ($z_t^{(2)}$) is supervised by boundary-aware many-hot targets to provide denser gradients. Fig. \ref{fig:head} shows the details of this head.

\begin{figure}[!t]
\centering
\resizebox{\columnwidth}{!}{%
\begin{tikzpicture}[
    font=\scriptsize,
    node distance=0.45cm and 0.75cm,
    box/.style={draw, rounded corners=2pt, thick, minimum height=0.58cm, text width=2.55cm, align=center},
    arr/.style={-{Stealth[length=4pt]}, thick}
]
\node[box, fill=blue!10] (feat) {Feature vector\\4864-D};
\node[box, fill=green!12, below=of feat] (fc) {Linear $4864\rightarrow1024$\\ReLU};
\node[box, fill=green!8, below=of fc] (drop) {Dropout\\$p=0.5$};
\node[box, fill=orange!18, below left=0.5cm and 0.25cm of drop] (one) {Primary branch\\one-hot target};
\node[box, fill=purple!14, below right=0.5cm and 0.25cm of drop] (many) {Auxiliary branch\\many-hot target};
\node[box, fill=red!8, below=1.45cm of drop, text width=5.35cm] (loss) {$\mathcal{L}_{\mathrm{total}}=\mathcal{L}_{\mathrm{focal}}^{(1)}+0.3\mathcal{L}_{\mathrm{focal}}^{(2)}$};
\draw[arr] (feat) -- (fc);
\draw[arr] (fc) -- (drop);
\draw[arr] (drop.south) -- ++(0,-0.18) -| (one.north);
\draw[arr] (drop.south) -- ++(0,-0.18) -| (many.north);
\draw[arr] (one.south) |- ([xshift=-1.3cm]loss.north);
\draw[arr] (many.south) |- ([xshift=1.3cm]loss.north);
\end{tikzpicture}
}
\caption{AutoShotV2 classification head. The primary branch preserves sharp boundary localization, while the auxiliary branch supplies denser supervision around each boundary.}
\label{fig:head}
\end{figure}

\subsection{Focal Loss for Class Imbalance}

Due to the extreme sparsity of boundary frames ($<1\%$), standard Binary Cross-Entropy (BCE) loss tends to be dominated by easy non-boundary frames, causing the model to bias toward predicting negative classes. We address this class imbalance by optimizing the classification head using Focal Loss \cite{lin2017focal}:
\begin{equation}
    \mathcal{L}_{\mathrm{focal}} = -\alpha(1-p_t^{*})^{\gamma}\log(p_t^{*}),
\end{equation}
where $p_t^* = p_t$ if the target $y_t=1$ and $1 - p_t$ otherwise, with $p_t = \sigma(z_t)$. The dynamic scaling factor $(1-p_t^*)^{\gamma}$ automatically down-weights the loss contributions of easy-to-classify frames and concentrates training gradients on hard boundaries and complex transitions. Grid search on the validation set yields the optimal parameters $\gamma=1.0$ and $\alpha=0.5$.

\subsection{Boundary-Aware Many-Hot Supervision}

One-hot annotations are temporally precise but sparse. They also do not account for labeling offsets or the transition span of gradual transitions. To address this, AutoShotV2 introduces an auxiliary many-hot supervision target $y_{\mathrm{mh}}[t]$ that expands the positive labels to adjacent frames, inspired by label smoothing techniques \cite{Muller2019LabelSmoothing}:
\begin{equation}
    y_{\mathrm{mh}}[t] =
    \begin{cases}
        1, & \exists\, t_b \;\mathrm{s.t.}\; |t-t_b|\leq1,\\
        0, & \mathrm{otherwise},
    \end{cases}
\end{equation}
where $t_b$ represents a ground-truth boundary frame.

The overall training objective combines both branches:
\begin{equation}
    \mathcal{L}_{\mathrm{total}} =
    \mathcal{L}_{\mathrm{focal}}(\hat{y}^{(1)},y_{\mathrm{oh}})
    +0.3\,\mathcal{L}_{\mathrm{focal}}(\hat{y}^{(2)},y_{\mathrm{mh}}).
\end{equation}
The auxiliary weight $0.3$ prevents the many-hot branch from over-smoothing the final boundary localization.

\subsection{Probability Calibration and Post-Processing}

\textbf{Temperature Scaling.} To calibrate the output confidence scores, the primary branch logit $z_t^{(1)}$ is scaled by a temperature parameter $T$:
\begin{equation}
    p_t^{\mathrm{cal}} = \sigma\left(\frac{z_t^{(1)}}{T}\right).
\end{equation}
The optimal temperature $T^*$ is found by minimizing the binary cross-entropy on the validation set:
\begin{equation}
    T^* = \arg\min_{T\in[0.1,10.0]}
    \mathcal{L}_{\mathrm{BCE}}\left(\sigma\left(\frac{z^{(1)}}{T}\right),y\right),
\end{equation}
resulting in $T^*=\PaperTemperature$. Because $T^* < 1$, temperature scaling sharpens the probability peaks for confident predictions.

\textbf{Gaussian Temporal Smoothing.} SBD models are prone to generating false-positive spikes due to single-frame disturbances (e.g., flash effects, camera shakes, quick occlusion). We apply a 1-D Gaussian filter to smooth the calibrated probability sequence:
\begin{equation}
    \tilde{p}_t = \sum_{k} G_{\sigma}(k)p_{t-k}^{\mathrm{cal}},
\end{equation}
where the Gaussian kernel is defined as:
\begin{equation}
    G_{\sigma}(k)= \frac{1}{\sqrt{2\pi}\sigma}\exp\left(-\frac{k^2}{2\sigma^2}\right).
\end{equation}
We select $\sigma=\PaperGaussianSigma$ via validation search, which suppresses isolated high-frequency false-alarm spikes while preserving temporally supported transition segments. EMA is not part of this deployment pipeline; it is evaluated separately in a full-model fine-tuning study and is therefore excluded from the controlled Phase 2 ablation.
